% TOP_PROBLEM: {"id": 9500001, "title": "Probabilistic McMillan theorem in higher dimensions", "category_id": 19, "category": {"id": 19, "name": "probability", "display_name": "Probability", "description": "Problems involving probability theory, stochastic processes, and random structures.", "slug": "probability", "order_index": 19, "created_at": "2026-07-31T15:26:25.671Z"}, "status": "partially_solved", "problem_number": "AMR-094-0001", "set_id": 13, "statusQualification": "For arbitrary unbounded domains the assertion is on the event of finite exit, where the source expression X_tau is defined."}
% FIRST_POSTED: 2026-10-01
% ENTRY_KIND: statement_only
% UnsolvedMath ID 9500001; source catalogue code AMR-094-0001
% !TeX program = lualatex
% Definition and sources only; this file is not an LLM solution attempt.
\documentclass[11pt,a4paper]{article}
\usepackage[margin=25mm]{geometry}
\usepackage{fontspec}
\setmainfont{lmroman10-regular.otf}[BoldFont=lmroman10-bold.otf,
 ItalicFont=lmroman10-italic.otf,BoldItalicFont=lmroman10-bolditalic.otf]
\usepackage{amsmath,amssymb,mathtools}
\usepackage{unicode-math}
\setmathfont{latinmodern-math.otf}
\AtBeginDocument{\let\setminus\smallsetminus}
\usepackage{xurl,hyperref}
\hypersetup{colorlinks=true,urlcolor=blue}
\setcounter{secnumdepth}{0}
\setlength{\parindent}{0pt}
\setlength{\parskip}{5pt}
\setlength{\emergencystretch}{3em}
\begin{document}
\section*{Probabilistic McMillan theorem in higher dimensions}\label{9500001}
{\small UnsolvedMath ID 9500001 · AMR-094-0001 · Probability · AMR Open Problem Lists}
\subsection{Definitions and mathematical statement}
Let $d\ge3$ and let $X=(X_t)_{t\ge0}$ be standard Brownian motion in $\mathbb R^d$, started at $0$. Thus its coordinates are independent real Brownian motions with variance $t$ at time $t$. Let $D\subseteq\mathbb R^d$ be open with $0\in D$, and define the exit time $\tau=\inf\{t>0:X_t\notin D\}$.

On the event $\tau<\infty$, continuity gives $X_\tau\in\partial D$. For $t<\tau$, put $u_t=(X_t-X_\tau)/|X_t-X_\tau|$, a point of the unit sphere $S^{d-1}$. Let $A$ consist of all limits of sequences $u_{t_n}$ with $t_n\uparrow\tau$. Equivalently,
\[
A=\bigcap_{0<s<\tau}\overline{\{u_t:s<t<\tau\}}\subseteq S^{d-1}.
\]
Must $A$, almost surely on finite exit, equal either the whole sphere or a closed hemisphere $\{u\in S^{d-1}:u\cdot v\ge0\}$ for some unit vector $v$?

The hemisphere and the choice between the two alternatives may depend on the sample path. The assertion is universal over domains, without a smoothness or connected-boundary assumption. The finite-exit qualification is necessary because an arbitrary unbounded domain need not be exited, and then $X_\tau$ is not defined. For bounded domains it is automatic. The question concerns the entire set of limiting approach directions, rather than the existence of a single tangent direction or a limiting normal vector.
\subsection{Short English statement}
At a finite Brownian exit from any higher-dimensional domain, are the limiting approach directions always a sphere or a closed hemisphere?
\subsection{Sources}
\begin{itemize}
\item[S1] ULAM.AI and the UnsolvedMath Contributors, supplied problems.json, record 9500001 (AMR-094-0001), AMR Open Problem Lists. Catalogue identity and source transcription; CC BY 4.0.
\url{https://huggingface.co/datasets/ulamai/UnsolvedMath}.
\item[S2] Burdzy - My favorite open problems (2009), Problem 1. Original source indexed by AMR-094-0001.
\url{https://sites.math.washington.edu/~burdzy/open_mathjax.php}.
\end{itemize}
\end{document}
